\subsection{Problem~2}%
\label{problem:2}

A balanced normal diet assumes that we should consume at least 60 grams of
proteins and at least 120 grams of carbohydrates daily. We assume that 100 grams of cheese
contains 20 grams of proteins and 20 grams of carbohydrates, but the same amount of bread
contains 10 grams of proteins and 30 grams of carbohydrates. Please determine the most
economical diet under the assumption that the price of cheese is 30\,PLN/kg and bread costs
20\,PLN/kg.

\subsubsection*{Solution}

Let $x_1$ (kg of cheese) and $x_2$ (kg of bread) be the decision variables.
Rescaling nutrition data to per-kg units ($1\,\text{kg} = 10 \times 100\,\text{g}$) gives the LP:

\begin{equation*}
\begin{split}
    \min_{x_1,x_2}\ & 30x_1 + 20x_2 \\
    \text{subject to}\quad
        & 200x_1 + 100x_2 \geq 60 \quad\text{(protein)} \\
        & 200x_1 + 300x_2 \geq 120 \quad\text{(carbohydrates)} \\
        & x_1,\, x_2 \geq 0
\end{split}
\end{equation*}

Because the constraints are of the $\geq$ form, the initial point $x_1 = x_2 = 0$ is
\emph{infeasible}: the basic solution obtained by adding slack variables would require
negative slacks.
Instead, surplus variables $s_1, s_2 \geq 0$ are subtracted and artificial variables
$a_1, a_2 \geq 0$ are added to obtain an initial basic feasible solution, which leads to
the two-phase simplex method~\cite{Zdunek}:

\begin{equation*}
    \mathbf{A} =
    \begin{bmatrix}
        200 & 100 & -1 & 0 \\
        200 & 300 &  0 & -1
    \end{bmatrix},\quad
    \mathbf{b} = \begin{bmatrix} 60 \\ 120 \end{bmatrix},\quad
    \mathbf{f} = \begin{bmatrix} 30 & 20 & 0 & 0 \end{bmatrix}
\end{equation*}

\paragraph{Phase~1.}
The auxiliary objective $\min\,(a_1 + a_2)$ is minimised over the extended system
$\mathbf{A}x + a = \mathbf{b}$.
The initial Phase~1 tableau (columns: $x_1,x_2,s_1,s_2,a_1,a_2,b$) is:

\begin{table}[H]
    \centering
    \begin{tabular}{c|cc cc cc|c}
         & $x_1$ & $x_2$ & $s_1$ & $s_2$ & $a_1$ & $a_2$ & $b$ \\
        \hline
        $a_1$ & 200 & 100 & $-1$ & 0 & 1 & 0 & 60 \\
        $a_2$ & 200 & 300 & 0 & $-1$ & 0 & 1 & 120 \\
        \hline
        $z_1$ & $-400$ & $-400$ & 1 & 1 & 0 & 0 & $-180$ \\
    \end{tabular}
\end{table}

After two pivot operations the Phase~1 objective reaches~0, confirming feasibility, and
the BFS $x_1 = 0.15$, $x_2 = 0.30$ is already in the basis.

\paragraph{Phase~2.}
The artificial columns are discarded and the original objective row
$\mathbf{f} = [30, 20, 0, 0]$ is substituted.
After eliminating the basis variables from the objective row, all reduced costs are
non-negative, so the algorithm terminates at the Phase~1 BFS.
The optimal solution is:

\begin{equation*}
    \systeme{x_1 = 0.15\,\text{kg (cheese)},x_2 = 0.30\,\text{kg (bread)}}
\end{equation*}

Both nutritional lower bounds are tight at this vertex:
\begin{equation*}
    200(0.15) + 100(0.30) = 60\,\text{g protein},\quad
    200(0.15) + 300(0.30) = 120\,\text{g carbohydrates.}
\end{equation*}
The minimum daily cost is $30(0.15) + 20(0.30) = \mathbf{10.50}\,\text{PLN}$.

\subsubsection*{Matlab Implementation}
\label{problem:2:impl}
\lstinputlisting[style=Matlab-editor, breaklines=false]{problems/Problem_2.m}
